% TOP_PROBLEM: {"id": 30006796, "title": "Polynomial-Time Distribution-Free PAC Learning of DNF Formulas", "category_id": 15, "category": {"id": 15, "name": "computer_science", "display_name": "Computer Science", "description": "Computational complexity, algorithms, and theoretical CS.", "slug": "computer-science", "order_index": 15, "created_at": "2026-07-31T15:26:25.671Z"}, "status": "open"}
% FIRST_POSTED: 2026-09-27
% !TeX program = lualatex
% Compile twice: lualatex 30006796.tex
% All references and prose are embedded; no BibTeX database or external inputs.
\documentclass[11pt,a4paper]{article}
\usepackage[margin=24mm,headheight=14pt]{geometry}
\usepackage{fontspec}
\setmainfont{Latin Modern Roman}
\setsansfont{Latin Modern Sans}
\setmonofont{Latin Modern Mono}
\usepackage{amsmath,amssymb,mathtools}
\usepackage{unicode-math}
\setmathfont{Latin Modern Math}
\usepackage{microtype}
\usepackage{enumitem,xurl,needspace}
\usepackage[dvipsnames]{xcolor}
\usepackage{hyperref}
\hypersetup{unicode=true,colorlinks=true,linkcolor=MidnightBlue,urlcolor=MidnightBlue,
 pdftitle={Research Notebook: Section 30006796},pdfauthor={Research notebook},bookmarksnumbered=true}
\usepackage{bookmark}
\usepackage{tocloft}
\setlength{\cftsecnumwidth}{3em}
\cftsetpnumwidth{2.5em}
\cftsetrmarg{4em}
\usepackage{titlesec}
\titleformat{\section}{\Large\bfseries\raggedright}{\thesection}{0.7em}{}
\usepackage{fancyhdr}
\pagestyle{fancy}\fancyhf{}
\fancyhead[L]{\small\sffamily Open Problems | Research notebook}
\fancyhead[R]{\small 219 | 27 September 2026}
\fancyfoot[C]{\thepage}
\setlength{\parindent}{0pt}
\setlength{\parskip}{5pt plus 1pt}
\setlength{\emergencystretch}{3em}
\setcounter{tocdepth}{1}
\setcounter{secnumdepth}{1}
\setlist[itemize]{leftmargin=3em,itemsep=5pt,topsep=4pt,parsep=0pt}
\allowdisplaybreaks
\newcommand{\N}{\mathbb{N}}
\newcommand{\Z}{\mathbb{Z}}
\newcommand{\Q}{\mathbb{Q}}
\newcommand{\R}{\mathbb{R}}
\newcommand{\C}{\mathbb{C}}
\newcommand{\F}{\mathbb{F}}
\newcommand{\A}{\mathbb{A}}
\newcommand{\PP}{\mathbb P}
\newcommand{\E}{\mathbb{E}}
\newcommand{\eps}{\varepsilon}
\newcommand{\dd}{\,\mathrm d}
\newcommand{\sref}[2]{\hyperlink{source:#1:#2}{[S#2]}}
\newcommand{\eref}[2]{\hyperlink{extra:#1:#2}{[E#2]}}
% Turn the next switch off for a shorter reading copy. The quotations preserve
% every exactTarget field, including qualifications omitted from a short summary.
\newif\ifshowcatalogue
\showcataloguetrue
\newcommand{\cataloguescope}[1]{\ifshowcatalogue
 \paragraph{Exact scope in the supplied catalogue.}
 \begin{quote}\small #1\end{quote}\fi}

% Independently compilable extraction; original section text retained.
\numberwithin{equation}{section}
\begin{document}
\setcounter{section}{0}
% ================================================================
% problemId: 30006796
% Canonical problem ID: 30006796
% exactTarget SHA-256: a2c4b2de3c6cc38835b336cb611743c1863c7928f8b10bdab2b7063d77f85d6f
\clearpage
\section*{\texorpdfstring{Polynomial-Time Distribution-Free PAC Learning of DNF Formulas}{Polynomial-Time Distribution-Free PAC Learning of DNF Formulas}}\label{30006796}
\subsection{Definitions and mathematical statement}
An $s$-term DNF on $\{0,1\}^n$ is an OR of at most $s$ terms, each an AND of literals $x_i$ or $\neg x_i$. The learner receives independent labeled examples $(x,f(x))$ with $x\sim D$, where $D$ is arbitrary and unknown. For $0<\eps,\delta<1$, it must output an efficiently evaluable hypothesis $h$ such that
\[
 \Pr_{\mathrm{samples,coins}}\bigl[\Pr_{x\sim D}(h(x)\ne f(x))\le\eps\bigr]\ge1-\delta
\]
using time polynomial in $n,s,1/\eps,\log(1/\delta)$. The hypothesis need not itself be a DNF. Membership queries or a known special input distribution are not supplied by the target.
\par\smallskip\textit{Formulation and concept references:} \sref{30006796}{1}.
\cataloguescope{Given an unknown s-term DNF f on \{0,1\}\textasciicircum{}n and random labeled examples drawn from an arbitrary unknown distribution D, construct a possibly improper hypothesis in time polynomial in n, s, 1/epsilon, and log(1/delta) whose D-error is at most epsilon with probability at least 1-delta.}
\subsection{Short English statement}
Can arbitrary DNF formulas be learned efficiently from random labeled examples drawn from any unknown distribution, even when the learner may output a different kind of hypothesis?
% BEGIN RESEARCH ATTEMPT 219
\subsection{Research attempt}

\paragraph{Current frontier.}
The 2025 paper supplied with the entry develops a quasipolynomial-time
procedure with \emph{membership queries}, including results for DNFs whose
terms have equal length. Its Section~1.1 explicitly separates these models
from distribution-free example-only learning. Its locally mixing walk
therefore cannot simply be run using the sample oracle in this target.
\sref{30006796}{1}, Abstract and Section~1.1.
A separate line of hardness results constrains learners that must output
DNF hypotheses of specified sizes; that is not automatically a barrier to
the unrestricted improper learner requested here. \eref{30006796}{1}, Abstract
and Introduction. No full example-only polynomial learner is obtained from
these sources. The attempt below isolates an algorithmic task that would
suffice, and tests an overly strong implementation of it.

\paragraph{Attack route.}
Exhaustive empirical risk minimization has enough samples but too much
computation. A random-walk strategy needs labels at points the unknown
distribution may never sample. A third route greedily covers positive
examples with conjunctions that exclude all negative examples. The
pigeonhole principle supplies a useful term at every stage. The missing
step is to \emph{find} one at polynomial cost, without enumerating all
$3^n$ conjunctions. This route may output a larger DNF, so its sample
bound must be checked for that larger hypothesis class, not for the
unknown $s$-term target alone.

\paragraph{Attempt.}
\textbf{1. Sample complexity is not the obstruction.}
Allow each variable in a term to appear positively, negatively, or not at
all. There are at most $3^n$ noncontradictory terms. Padding lists by an
empty choice bounds the number of $s$-term DNFs by $(3^n+1)^s$.
For any hypothesis with true error greater than $\eps$, the probability
of agreeing with $m$ independent realizable examples is at most
$(1-\eps)^m\le e^{-m\eps}$. A union bound gives a consistent learner
using
\[
 m\ge\frac{s\log(3^n+1)+\log(1/\delta)}{\eps}.
\]
Searching the entire hypothesis class meets this statistical guarantee,
but costs exponentially many tests. This calculation neither bounds that
search by a polynomial nor establishes a computational lower bound.

\textbf{2. A precise conditional greedy learner.}
Consider the following promised finite-sample search task, denoted
\emph{weak consistent-term search}. Given a labeled sample realizable by
an $s$-term DNF, a subset $R$ of its positive examples, and the full list
of its negative examples, return a conjunction $T$ which accepts no
listed negative and accepts at least $|R|/s$ of the examples in $R$.
Examples are counted with multiplicity. A valid conjunction always
exists: each target term excludes every true negative, and the target's
at most $s$ terms cover $R$.

Suppose a uniform algorithm solves this task in time polynomial in
$n,s,m$. Start with all positive examples in $R$, find a term, delete
its covered positives, and repeat. With $r$ terms selected, the number
left is at most $m(1-1/s)^r\le me^{-r/s}$ for $s>1$.
The $s=1$ case finishes in one step. In either case at most
\[
 T_m=\lceil s\log m\rceil+1\le s\log m+2
\]
terms suffice, with the natural one-step convention if $m=1$.
Their OR is consistent with the whole sample and is evaluable in
$O(nT_m)$ operations.

Choose $m\ge2$ by doubling until
\begin{equation}\label{30006796:sample}
 m\eps\ge n(s\log m+2)\log4+\log(1/\delta).
\end{equation}
This loop terminates with $m$ bounded by a polynomial in
$n,s,1/\eps,\log(1/\delta)$: for example, putting
$B=ns+\log(1/\delta)+2$, an absolute sufficiently large constant times
$(B/\eps)^2$ satisfies the inequality, using $\log u\le\sqrt u$ for
large $u$. The class of all DNFs of at most $T_m$ terms has cardinality
at most $4^{nT_m}$. Applying the same union bound to that fixed class,
\eqref{30006796:sample} proves that \emph{every} consistent output of the
procedure has error at most $\eps$ with probability at least $1-\delta$.
The number of search calls is also polynomial.

Thus a polynomial-time weak consistent-term search algorithm would prove
the full selected learning assertion. The proof handles arbitrary $D$
and uses examples only. It does not assume that the hypothesis has just
$s$ terms. It also does not prove the converse: an arbitrary improper
PAC learner need not solve this particular finite-sample search task.

\textbf{3. Exact optimization is too strong a proposed implementation.}
An apparent implementation is to find a consistent term covering the
\emph{largest} number of positive examples. The following reduction
shows a concrete difficulty even for a term required to accept the
all-one point.

Let $G$ be a simple graph on $\{1,\ldots,n\}$. Use an anchor
$a=(1,\ldots,1)$, and for each vertex $i$ a positive example $p^i$ which
is zero at coordinate $i$ and one elsewhere. For each edge $\{i,j\}$,
use a negative example $q^{ij}$ which is zero exactly at $i,j$.
A conjunction accepting $a$ has only positive literals, indexed by a
subset $S$. It rejects every $q^{ij}$ exactly when $S$ meets every edge,
that is, when $S$ is a vertex cover. It accepts $p^i$ exactly when
$i\notin S$. Consequently
\begin{equation}\label{30006796:independent}
 \max_T|\{i:T(p^i)=1\}|=\alpha(G),
\end{equation}
where the maximum is over consistent conjunctions accepting $a$ and
$\alpha(G)$ is the maximum independent-set size.
To remove the explicit anchor requirement for an exact optimizer, give
$a$ more than $n$ positive copies. Any conjunction missing $a$ loses
more positives than it could gain among all $p^i$; a consistent term
containing the anchor exists, namely the conjunction of all variables.
Thus an exact maximum-coverage optimizer yields a maximum independent
set in polynomial time.

This is an exact reduction, not a claim that the \emph{weak} $1/s$
search task is NP-hard. Indeed the constructed sample may need many
terms to realize all its positives, making the weak guarantee small.
Nor does \eqref{30006796:independent} show that a general improper PAC
learner decides NP-complete problems. The familiar hardness of maximum
independent set is used only to reject this stronger proposed optimizer.
\eref{30006796}{2}, the clique/vertex-cover reductions.

\textbf{4. Two tests of distribution-dependent shortcuts.}
First, a long conjunction need not be rare. If $D$ is concentrated at
the all-one point, the conjunction $x_1\wedge\cdots\wedge x_n$ has
probability one. Deleting long terms by the uniform-distribution bound
$2^{-\text{length}}$ therefore fails at the first step.

Second, suppose $f\equiv1$ and $D$ is uniform on the $n$ standard basis
vectors $e_1,\ldots,e_n$. The most specific conjunction accepting
$k\ge2$ \emph{distinct} observed positives contains negative literals
on every as-yet-unobserved coordinate. It accepts exactly the $k$
observed basis vectors, and has mass $k/n$, although the empty target
term accepts everything. Thus positive consensus from a small seed does
not automatically find a term covering a constant fraction of the
positive distribution. This example does not rule out algorithms that
successfully generalize beyond the most specific conjunction.

The companion finite computation checks \eqref{30006796:independent} by
exhaustive subset enumeration on several small graphs. Its purpose is to
verify the reduction, not to infer an asymptotic complexity separation.

\paragraph{Stress test.}
The sample bound is applied to the \emph{output} class of $T_m$-term
DNFs; applying it only to $s$ terms would be unjustified. Every selected
term is consistent with the original negative sample, rather than merely
with negatives retained after deletions. The search promise is preserved
because the same target DNF covers every residual positive subset.
The proposed oracle is uniform and polynomial in the sample size; it is
not implemented by $3^n$ enumeration under another name. No membership
query or assumption on marginal probabilities appears in the reduction.

There is no inference of $\mathrm P=\mathrm{NP}$ from the target. The
maximum-coverage subroutine is stronger than necessary, and the proper
learning limitations in \eref{30006796}{1} have an output restriction absent
from the question. Those distinctions are substantive, not terminological.

\paragraph{Outcome.}
\textbf{Outcome: Conditional reduction.}

A polynomial weak consistent-term search procedure would give an
example-only distribution-free learner, with polynomial running time,
efficient evaluation, and the required dependence on $\log(1/\delta)$.
The full reduction and its generalization estimate are supplied here.
An exact-optimization implementation encounters a maximum-independent-set
reduction, and two distributional heuristics have explicit counterexamples.

The weak search procedure itself remains unconstructed. Therefore this
is not a polynomial-time DNF learner, and it does not exclude a different
improper learner. No priority claim is made for the elementary counting
or greedy-cover principles used in the argument.
\needspace{20\baselineskip}
% END RESEARCH ATTEMPT 219
\subsection{Sources}
\begin{itemize}\raggedright
\item[\textnormal{[S1]}]\hypertarget{source:30006796:1}{} Alman, Nadimpalli, Patel, and Servedio, DNF Learning via Locally Mixing Random Walks, Section 1.1.
\url{https://arxiv.org/html/2505.18839v1}.
{\small\textit{Catalogue formulation source.}}
% BEGIN ADDED REFERENCES 219
\item[\textnormal{[E1]}]\hypertarget{extra:30006796:1}{}
N. H. Bshouty, \emph{Superpolynomial Lower Bounds for Learning Monotone
Classes}, arXiv:2301.08486, Abstract and Introduction (the DNF-output
and complexity-assumption restrictions are retained).
\url{https://arxiv.org/abs/2301.08486}. Consulted 27 September 2026.
\item[\textnormal{[E2]}]\hypertarget{extra:30006796:2}{}
R. M. Karp, \emph{Reducibility Among Combinatorial Problems}, in
\emph{Complexity of Computer Computations}, Plenum, 1972, 85--103,
clique and node-cover problems and their polynomial reductions.
\url{https://doi.org/10.1007/978-1-4684-2001-2_9}.
Used only for the complexity interpretation of the explicitly proved
finite reduction, not a lower bound for improper PAC learning.
% END ADDED REFERENCES 219
\end{itemize}

\end{document}
